\subsection{Interação com o manipulador robótico em base livre}

Ao iniciar a fase de proximidade, o movimento do manipulador influencia a atitude e a posição do \emph{chaser}, devido ao acoplamento dinâmico entre a base e o MR.  
A movimentação das juntas gera torques internos e estes se propagam para a base.  
Sendo assim, mesmo leves variações de ângulo nas juntas podem gerar contribuições na atitude da base.
Durante as manobras de aproximação curta, o controlador de atitude compensa essas reações inerciais e coordena os torques do MR e da base.  